% Options for packages loaded elsewhere
% Options for packages loaded elsewhere
\PassOptionsToPackage{unicode}{hyperref}
\PassOptionsToPackage{hyphens}{url}
\PassOptionsToPackage{dvipsnames,svgnames,x11names}{xcolor}
%
\documentclass[
  letterpaper,
  DIV=11,
  numbers=noendperiod]{scrartcl}
\usepackage{xcolor}
\usepackage{amsmath,amssymb}
\setcounter{secnumdepth}{5}
\usepackage{iftex}
\ifPDFTeX
  \usepackage[T1]{fontenc}
  \usepackage[utf8]{inputenc}
  \usepackage{textcomp} % provide euro and other symbols
\else % if luatex or xetex
  \usepackage{unicode-math} % this also loads fontspec
  \defaultfontfeatures{Scale=MatchLowercase}
  \defaultfontfeatures[\rmfamily]{Ligatures=TeX,Scale=1}
\fi
\usepackage{lmodern}
\ifPDFTeX\else
  % xetex/luatex font selection
\fi
% Use upquote if available, for straight quotes in verbatim environments
\IfFileExists{upquote.sty}{\usepackage{upquote}}{}
\IfFileExists{microtype.sty}{% use microtype if available
  \usepackage[]{microtype}
  \UseMicrotypeSet[protrusion]{basicmath} % disable protrusion for tt fonts
}{}
\makeatletter
\@ifundefined{KOMAClassName}{% if non-KOMA class
  \IfFileExists{parskip.sty}{%
    \usepackage{parskip}
  }{% else
    \setlength{\parindent}{0pt}
    \setlength{\parskip}{6pt plus 2pt minus 1pt}}
}{% if KOMA class
  \KOMAoptions{parskip=half}}
\makeatother
% Make \paragraph and \subparagraph free-standing
\makeatletter
\ifx\paragraph\undefined\else
  \let\oldparagraph\paragraph
  \renewcommand{\paragraph}{
    \@ifstar
      \xxxParagraphStar
      \xxxParagraphNoStar
  }
  \newcommand{\xxxParagraphStar}[1]{\oldparagraph*{#1}\mbox{}}
  \newcommand{\xxxParagraphNoStar}[1]{\oldparagraph{#1}\mbox{}}
\fi
\ifx\subparagraph\undefined\else
  \let\oldsubparagraph\subparagraph
  \renewcommand{\subparagraph}{
    \@ifstar
      \xxxSubParagraphStar
      \xxxSubParagraphNoStar
  }
  \newcommand{\xxxSubParagraphStar}[1]{\oldsubparagraph*{#1}\mbox{}}
  \newcommand{\xxxSubParagraphNoStar}[1]{\oldsubparagraph{#1}\mbox{}}
\fi
\makeatother


\usepackage{longtable,booktabs,array}
\usepackage{calc} % for calculating minipage widths
% Correct order of tables after \paragraph or \subparagraph
\usepackage{etoolbox}
\makeatletter
\patchcmd\longtable{\par}{\if@noskipsec\mbox{}\fi\par}{}{}
\makeatother
% Allow footnotes in longtable head/foot
\IfFileExists{footnotehyper.sty}{\usepackage{footnotehyper}}{\usepackage{footnote}}
\makesavenoteenv{longtable}
\usepackage{graphicx}
\makeatletter
\newsavebox\pandoc@box
\newcommand*\pandocbounded[1]{% scales image to fit in text height/width
  \sbox\pandoc@box{#1}%
  \Gscale@div\@tempa{\textheight}{\dimexpr\ht\pandoc@box+\dp\pandoc@box\relax}%
  \Gscale@div\@tempb{\linewidth}{\wd\pandoc@box}%
  \ifdim\@tempb\p@<\@tempa\p@\let\@tempa\@tempb\fi% select the smaller of both
  \ifdim\@tempa\p@<\p@\scalebox{\@tempa}{\usebox\pandoc@box}%
  \else\usebox{\pandoc@box}%
  \fi%
}
% Set default figure placement to htbp
\def\fps@figure{htbp}
\makeatother





\setlength{\emergencystretch}{3em} % prevent overfull lines

\providecommand{\tightlist}{%
  \setlength{\itemsep}{0pt}\setlength{\parskip}{0pt}}



 


\KOMAoption{captions}{tableheading}
\makeatletter
\@ifpackageloaded{tcolorbox}{}{\usepackage[skins,breakable]{tcolorbox}}
\@ifpackageloaded{fontawesome5}{}{\usepackage{fontawesome5}}
\definecolor{quarto-callout-color}{HTML}{909090}
\definecolor{quarto-callout-note-color}{HTML}{0758E5}
\definecolor{quarto-callout-important-color}{HTML}{CC1914}
\definecolor{quarto-callout-warning-color}{HTML}{EB9113}
\definecolor{quarto-callout-tip-color}{HTML}{00A047}
\definecolor{quarto-callout-caution-color}{HTML}{FC5300}
\definecolor{quarto-callout-color-frame}{HTML}{acacac}
\definecolor{quarto-callout-note-color-frame}{HTML}{4582ec}
\definecolor{quarto-callout-important-color-frame}{HTML}{d9534f}
\definecolor{quarto-callout-warning-color-frame}{HTML}{f0ad4e}
\definecolor{quarto-callout-tip-color-frame}{HTML}{02b875}
\definecolor{quarto-callout-caution-color-frame}{HTML}{fd7e14}
\makeatother
\makeatletter
\@ifpackageloaded{caption}{}{\usepackage{caption}}
\AtBeginDocument{%
\ifdefined\contentsname
  \renewcommand*\contentsname{Table of contents}
\else
  \newcommand\contentsname{Table of contents}
\fi
\ifdefined\listfigurename
  \renewcommand*\listfigurename{List of Figures}
\else
  \newcommand\listfigurename{List of Figures}
\fi
\ifdefined\listtablename
  \renewcommand*\listtablename{List of Tables}
\else
  \newcommand\listtablename{List of Tables}
\fi
\ifdefined\figurename
  \renewcommand*\figurename{Figure}
\else
  \newcommand\figurename{Figure}
\fi
\ifdefined\tablename
  \renewcommand*\tablename{Table}
\else
  \newcommand\tablename{Table}
\fi
}
\@ifpackageloaded{float}{}{\usepackage{float}}
\floatstyle{ruled}
\@ifundefined{c@chapter}{\newfloat{codelisting}{h}{lop}}{\newfloat{codelisting}{h}{lop}[chapter]}
\floatname{codelisting}{Listing}
\newcommand*\listoflistings{\listof{codelisting}{List of Listings}}
\makeatother
\makeatletter
\makeatother
\makeatletter
\@ifpackageloaded{caption}{}{\usepackage{caption}}
\@ifpackageloaded{subcaption}{}{\usepackage{subcaption}}
\makeatother
\usepackage{bookmark}
\IfFileExists{xurl.sty}{\usepackage{xurl}}{} % add URL line breaks if available
\urlstyle{same}
\hypersetup{
  pdftitle={Lecture 1},
  colorlinks=true,
  linkcolor={blue},
  filecolor={Maroon},
  citecolor={Blue},
  urlcolor={Blue},
  pdfcreator={LaTeX via pandoc}}


\title{Lecture 1}
\author{}
\date{}
\begin{document}
\maketitle

\renewcommand*\contentsname{Table of contents}
{
\hypersetup{linkcolor=}
\setcounter{tocdepth}{3}
\tableofcontents
}

\section{Lecture 1: title of lecture}\label{lecture-1-title-of-lecture}

\subsection{Learning Objectives}\label{learning-objectives}

\begin{itemize}
\tightlist
\item
  \emph{Add your learning objectives here.}
\end{itemize}

\subsection{1. first section}\label{first-section}

\subsubsection{1.1. first subsection}\label{first-subsection}

\paragraph{1.1.1. first subsubsection}\label{first-subsubsection}

\begin{tcolorbox}[enhanced jigsaw, breakable, title=\textcolor{quarto-callout-note-color}{\faInfo}\hspace{0.5em}{Worked Example}, opacityback=0, coltitle=black, titlerule=0mm, opacitybacktitle=0.6, colback=white, bottomtitle=1mm, colframe=quarto-callout-note-color-frame, leftrule=.75mm, bottomrule=.15mm, toptitle=1mm, arc=.35mm, rightrule=.15mm, toprule=.15mm, colbacktitle=quarto-callout-note-color!10!white, left=2mm]

\textbf{Example 1:}\\
Suppose you invest \$1,000 in a bond earning 5\% interest per year for 3
years.\\
- What is the future value?\\
\[
FV = 1000 \times (1 + 0.05)^3 = 1000 \times 1.157625 = \$1,157.63
\]

\textbf{Explanation:}\\
We apply the compound interest formula for three years.

\end{tcolorbox}

\begin{tcolorbox}[enhanced jigsaw, breakable, title=\textcolor{quarto-callout-tip-color}{\faLightbulb}\hspace{0.5em}{Quick Application}, opacityback=0, coltitle=black, titlerule=0mm, opacitybacktitle=0.6, colback=white, bottomtitle=1mm, colframe=quarto-callout-tip-color-frame, leftrule=.75mm, bottomrule=.15mm, toptitle=1mm, arc=.35mm, rightrule=.15mm, toprule=.15mm, colbacktitle=quarto-callout-tip-color!10!white, left=2mm]

If the interest rate doubles to 10\% but the period is only 1 year,\\
- New future value calculation: \[
FV = 1000 \times (1+0.10)^1 = \$1,100
\]

\end{tcolorbox}

\begin{tcolorbox}[enhanced jigsaw, breakable, title=\textcolor{quarto-callout-important-color}{\faExclamation}\hspace{0.5em}{Try Yourself}, opacityback=0, coltitle=black, titlerule=0mm, opacitybacktitle=0.6, colback=white, bottomtitle=1mm, colframe=quarto-callout-important-color-frame, leftrule=.75mm, bottomrule=.15mm, toptitle=1mm, arc=.35mm, rightrule=.15mm, toprule=.15mm, colbacktitle=quarto-callout-important-color!10!white, left=2mm]

\textbf{Practice:}\\
Calculate the present value of \$2,000 to be received in 4 years if the
discount rate is 6\%.\\
(Hint: Use ( PV = \frac{F}{(1+r)^n} ))

\end{tcolorbox}

\subsection{Main Concepts}\label{main-concepts}

\begin{itemize}
\tightlist
\item
  \emph{Main concept 1}
\item
  \emph{Main concept 2}
\item
  \emph{Main concept 3}
\end{itemize}

\subsection{Definitions}\label{definitions}

\textbf{Term 1.} \emph{Definition placeholder.}

\textbf{Term 2.} \emph{Definition placeholder.}

\subsection{Key Formulas}\label{key-formulas}

\[
\text{Formula placeholder: } PV = \frac{F}{(1+r)^n}
\]

\subsection{Suggested Readings}\label{suggested-readings}

\begin{itemize}
\tightlist
\item
  \emph{Textbook chapter or paper reference}
\item
  \emph{Supplementary material}
\end{itemize}




\end{document}
